\documentclass[10pt,twoside]{article}\usepackage[]{graphicx}\usepackage[dvipsnames,svgnames,table]{xcolor}
% maxwidth is the original width if it is less than linewidth
% otherwise use linewidth (to make sure the graphics do not exceed the margin)
\makeatletter
\def\maxwidth{ %
  \ifdim\Gin@nat@width>\linewidth
    \linewidth
  \else
    \Gin@nat@width
  \fi
}
\makeatother

\definecolor{fgcolor}{rgb}{0.345, 0.345, 0.345}
\newcommand{\hlnum}[1]{\textcolor[rgb]{0.686,0.059,0.569}{#1}}%
\newcommand{\hlsng}[1]{\textcolor[rgb]{0.192,0.494,0.8}{#1}}%
\newcommand{\hlcom}[1]{\textcolor[rgb]{0.678,0.584,0.686}{\textit{#1}}}%
\newcommand{\hlopt}[1]{\textcolor[rgb]{0,0,0}{#1}}%
\newcommand{\hldef}[1]{\textcolor[rgb]{0.345,0.345,0.345}{#1}}%
\newcommand{\hlkwa}[1]{\textcolor[rgb]{0.161,0.373,0.58}{\textbf{#1}}}%
\newcommand{\hlkwb}[1]{\textcolor[rgb]{0.69,0.353,0.396}{#1}}%
\newcommand{\hlkwc}[1]{\textcolor[rgb]{0.333,0.667,0.333}{#1}}%
\newcommand{\hlkwd}[1]{\textcolor[rgb]{0.737,0.353,0.396}{\textbf{#1}}}%
\let\hlipl\hlkwb

\usepackage{framed}
\makeatletter
\newenvironment{kframe}{%
 \def\at@end@of@kframe{}%
 \ifinner\ifhmode%
  \def\at@end@of@kframe{\end{minipage}}%
  \begin{minipage}{\columnwidth}%
 \fi\fi%
 \def\FrameCommand##1{\hskip\@totalleftmargin \hskip-\fboxsep
 \colorbox{shadecolor}{##1}\hskip-\fboxsep
     % There is no \\@totalrightmargin, so:
     \hskip-\linewidth \hskip-\@totalleftmargin \hskip\columnwidth}%
 \MakeFramed {\advance\hsize-\width
   \@totalleftmargin\z@ \linewidth\hsize
   \@setminipage}}%
 {\par\unskip\endMakeFramed%
 \at@end@of@kframe}
\makeatother

\definecolor{shadecolor}{rgb}{.97, .97, .97}
\definecolor{messagecolor}{rgb}{0, 0, 0}
\definecolor{warningcolor}{rgb}{1, 0, 1}
\definecolor{errorcolor}{rgb}{1, 0, 0}
\newenvironment{knitrout}{}{} % an empty environment to be redefined in TeX

\usepackage{alltt}
\usepackage[marginparsep=1em]{geometry}
\geometry{lmargin=1.0in,rmargin=1.0in, bmargin=1.2in,  tmargin=1.2in}
\usepackage[dvipsnames,svgnames,table]{xcolor}
\usepackage{graphicx}
\usepackage{amssymb}
\usepackage{epstopdf}
\usepackage{verbatim}
\usepackage{enumerate}
\usepackage{bm}
\usepackage{amsthm}
\usepackage{float}
\usepackage{amsmath}
\usepackage{fancyhdr}
\usepackage{hyperref}
\usepackage{mathtools}

%%%%  SHORTCUT COMMANDS  %%%%
\newcommand{\ds}{\displaystyle}
\newcommand{\Z}{\mathbb{Z}}
\newcommand{\T}{\mathcal{T}}
\newcommand{\arc}{\rightarrow}
\newcommand{\R}{\mathbb{R}}
\newcommand{\RP}{\mathbb{R}(+)}
\newcommand{\Rs}{\mathbb{R}^{**}}
\newcommand{\C}{\mathbb{C}}
\newcommand{\E}{\mathbb{E}}
\newcommand{\B}{\mathcal{B}}
\newcommand{\PX}{\mathcal{P}(X)}
\newcommand*\rot{\rotatebox{90}}
\newcommand*\OK{\ding{51}}
\newcommand{\N}{\mathbb{N}}
\newcommand{\Q}{\mathbb{Q}}
\newcommand{\Answer}{\vspace{2mm}\textbf{\underline{Answer}}\\}
\newcolumntype{L}{>{$}l<{$}}
\newcolumntype{C}{>{$}c<{$}}
\newcommand{\GL}{\text{GL}}
\newcommand{\SL}{\text{SL}}

\newcommand{\stirling}[2]{\genfrac{\{}{\}}{0pt}{}{#1}{#2}}

\pagestyle{fancy}
\fancyhf{}
\renewcommand{\sectionmark}[1]{\markright{\thesection.\ #1}}
\lhead{\fancyplain{}{}}
\fancyhead[RE,RO]{Math 4451, Spring 2026}
\fancyfoot[RE,RO]{\thepage}
\IfFileExists{upquote.sty}{\usepackage{upquote}}{}
\begin{document}
\begin{flushright}
\begin{minipage}{.33\textwidth}
\rightline{YOUR NAME}
\rightline{\href{mailto:YOUR EMAIL}{YOUR EMAIL}}
\rightline{\today}
\end{minipage}
\end{flushright}

\begin{center}
{\large{\textbf{Homework 6}}}
\end{center}

\noindent\textbf{Note on Regularity Conditions:} For all problems in this assignment, you may assume that the standard Fisher regularity conditions hold.
Specifically, you may assume that the support of the distribution does not depend on the parameter $\theta$, and that the density functions are sufficiently smooth to allow the interchange of integration (or summation) and differentiation.
You may proceed directly to using the second derivative formulation:
$$
\mathcal{I}(\theta) = -E\left[ \frac{\partial^2}{\partial \theta^2} \ln f(X; \theta) \right]
$$

\begin{enumerate}

  \item (5 points) Let $X_1,\ldots, X_n$ be iid random variables with density function
  $$
  f(x_i; \theta) = (\theta + 1)x^\theta, \quad 0 \leq x \leq 1.
  $$

  \begin{enumerate}
    \item (2 points) Find the MLE of $\theta$.
    \item (3 points) Find the asymptotic variance of the MLE.
  \end{enumerate}


  \item (5 points) Let $X_1, \ldots, X_n$ be an iid sample from a Poisson distribution with parameter $\lambda > 0$. The pmf of a Poisson$(\lambda)$ random variable is given by:
  $$
  P(X = x) = \frac{\lambda^xe^{-\lambda}}{x!}, \quad \text{for } x \in \{0, 1, 2, \ldots\}
  $$

\begin{enumerate}
  \item (1 point) Calculate the Observed Fisher Information of all points, $I_n(\lambda)$.
  \item (1 points) Calculate the expected Fisher Information for all observations, $\mathcal{I}_n(\lambda)$.

  \item (1 point) Compare both $I_n(\lambda)$ and $\mathcal{I}_n(\lambda)$ at the MLE: $\lambda = \hat{\lambda}$.
  \item (2 points) Derive the asymptotic variance of the MLE, and compare to the exact variance of the estimator.
\end{enumerate}

\item (5 points) Let $X_1, \ldots, X_n$ be an iid sample from a Normal distribution with known mean $\mu = 0$, and unknown variance $\theta = \sigma^2$. The density of a single observation is:
$$
f(x; \theta) = \frac{1}{\sqrt{2\pi\theta}} \exp\Big(-\frac{x^2}{2\theta}\Big).
$$

\begin{enumerate}
  \item (2 point) Find the score function for $\theta$.
  \item (3 points) Calculate the (expected) Fisher Information $\mathcal{I}_n(\theta)$.
\end{enumerate}

\item (5 points) In class, it was briefly argued that if $X_1, \ldots, X_n$ are iid, then the joint information $\mathcal{I}(\theta)$ is just $n$ times the individual information:
$$
\mathcal{I}(\theta) = n\mathcal{I}_1(\theta).
$$
Your task is to generalize this result for independent (not necessarily identical) random variables.

\begin{itemize}
  \item Let $X$ and $Y$ be independent random variables with density functions $f_X(x; \theta)$ and $f_Y(y; \theta)$, respectively (note that they do share the same parameter value $\theta$, but do not necessarily have to come from the same distribution).
  Show that if $\mathcal{I}_X(\theta)$, $\mathcal{I}_Y(\theta)$, and $\mathcal{I}_{X, Y}(\theta)$ are the individual informations and joint information (based on the joint density), respectively, then:
  $$
  \mathcal{I}_{X, Y}(\theta) = \mathcal{I}_X(\theta) + \mathcal{I}_Y(\theta)
  $$
\end{itemize}

\end{enumerate}


\end{document}
